\section{Tổng quan đề tài}

\subsection{Động cơ của đề tài}
    Trong bối cảnh toàn cầu thúc đẩy phát triển bền vững, việc chuyển dịch sang các phương tiện giao thông xanh - đặc biệt là xe điện - đang trở thành xu thế tất yếu. Tại Việt Nam, số lượng xe điện tăng trưởng nhanh chóng, nhưng đi kèm với đó là thách thức lớn về hạ tầng sạc. Người dùng xe điện không chỉ cần biết ở đâu có trạm sạc gần mình, mà còn mong muốn đảm bảo chắc chắn có chỗ sạc, theo dõi tiến trình sạc thuận tiện và thanh toán minh bạch. Khi những nhu cầu cơ bản này chưa được đáp ứng đầy đủ, trải nghiệm sử dụng xe điện trở nên hạn chế, làm chậm quá trình phổ cập phương tiện xanh.
\par
    Ở chiều ngược lại, các chủ trạm sạc - những đơn vị trực tiếp đầu tư hạ tầng - cũng gặp nhiều trở ngại. Họ cần một công cụ quản lý thống nhất để kiểm soát hoạt động của trạm, đánh giá hiệu suất khai thác và tương tác với khách hàng một cách chuyên nghiệp. Nếu thiếu đi hệ thống hỗ trợ, việc vận hành sẽ khó đạt hiệu quả kinh tế, giảm động lực mở rộng và phát triển mạng lưới.
\par
    Đồng thời, cơ quan quản lý toàn hệ thống/chủ hệ thống cũng không thể dựa vào các kênh rời rạc. Họ cần một nền tảng tập trung để giám sát, điều phối, và đưa ra quyết định dựa trên dữ liệu. Chỉ khi có cái nhìn tổng thể, hệ thống mới đảm bảo tính minh bạch, công bằng, và đủ sức phục vụ chiến lược phát triển giao thông điện tại quy mô quốc gia.
\par 
    Chính từ những nhu cầu giao thoa này - của người dùng, của chủ trạm, và của cơ quan quản lý - việc nghiên cứu và xây dựng một hệ thống quản lý trạm sạc xe điện mang ý nghĩa không chỉ về mặt công nghệ, mà còn về xã hội và môi trường. Đây là giải pháp thiết yếu để khắc phục rào cản hạ tầng, tạo động lực thúc đẩy quá trình chuyển đổi xanh, nâng cao trải nghiệm người dùng và tối ưu hóa hiệu quả vận hành toàn hệ thống.

\subsection{Mục tiêu của đề tài}
    Mục tiêu cốt lõi của đồ tài là xây dựng hệ thống quản lý trạm sạc xe máy điện, trong đó trọng tâm nguồn lực được ưu tiên cho việc tối ưu hóa trải nghiệm của người dùng cuối. Hệ thống hướng tới cung cấp một quy trình dịch vụ khép kín và liền mạch: từ việc định vị trạm sạc, quản lý đặt lịch, cho đến khả năng giám sát thông số sạc theo thời gian thực và thanh toán tự động. Bên cạnh việc thiết lập nền tảng quản trị cơ bản cho đơn vị vận hành, đề tài tập trung giải quyết triệt để các rào cản tiếp cận trạm sạc trong thực tế của người lái xe, qua đó đóng góp một giải pháp phần mềm mang tính ứng dụng cao, thúc đẩy quá trình phổ cập phương tiện xanh tại Việt Nam.

\subsection{Phạm vi của đề tài}
    Nhằm tối ưu hóa chất lượng phần mềm trong khuôn khổ thời gian thực hiện, đề tài tập trung nguồn lực phát triển chuyên sâu ứng dụng di động và hệ thống Backend phục vụ trực tiếp người sử dụng xe máy điện. Để đảm bảo tính toàn vẹn của kiến trúc tổng thể, các phân hệ quản trị dành cho chủ trạm được giới hạn phạm vi ở mức độ cung cấp dữ liệu nền tảng và thiết lập cấu hình cơ bản (như quản lý trạm sạc, cổng sạc). Sự phân bổ này cho phép đồ án dành tối đa không gian để giải quyết các bài toán kỹ thuật phức tạp ở luồng người dùng cuối, qua đó đảm bảo tính ổn định và hiệu năng của toàn hệ thống.
    \par
    Đối với ứng dụng di động, hệ thống tập trung hoàn thiện các nghiệp vụ cốt lõi nhằm tối ưu hóa trải nghiệm người dùng, bao gồm: tìm kiếm trạm sạc qua bản đồ số, quản lý đặt lịch sạc và giám sát trạng thái sạc theo thời gian thực. Ngoài ra ứng dụng còn được tích hợp với hệ thống thanh toán trực tuyến và cơ chế thông báo tự động, qua đó hình thành một chu trình dịch vụ khép kín và liền mạch.
    \par
    Về mặt giới hạn nghiên cứu, do thiếu hụt thiết bị phần cứng thực tế, quá trình giao tiếp giữa hệ thống trung tâm và trạm sạc được thực hiện thông qua phần mềm mô phỏng (OCPP ChargePoint Simulator), được phát triển và tuân thủ nghiêm ngặt theo đặc tả của giao thức \cite{ocpp_spec}. Song song đó, quy trình thanh toán trực tuyến được tích hợp và kiểm thử trên môi trường giả lập Sandbox của đối tác cung cấp dịch vụ. 

\subsection{Ý nghĩa của đề tài}
\subsubsection{Ý nghĩa thực tiễn}
    Đề tài mang lại giá trị ứng dụng rõ ràng trong bối cảnh Việt Nam đang thúc đẩy quá trình chuyển đổi sang giao thông xanh. Việc xây dựng một ứng dụng trạm sạc hoàn thiện giúp người dùng xe máy điện dễ dàng tiếp cận dịch vụ sạc một cách thuận tiện, minh bạch và an toàn. Thông qua việc giải quyết bài toán cốt lõi từ góc nhìn của người dùng, hệ thống góp phần xóa bỏ những lo ngại về hạ tầng năng lượng, trực tiếp khuyến khích cộng đồng chuyển đổi sang sử dụng phương tiện điện.
\subsubsection{Ý nghĩa khoa học}
    Về mặt học thuật, đề tài là một nghiên cứu tích hợp đa lĩnh vực, bao gồm phát triển hệ thống phần mềm, quản lý dữ liệu thời gian thực, xử lý tương tranh và thiết kế giao diện thân thiện với người dùng. Đề tài làm rõ cách tổ chức triển khai một hệ thống hướng người dùng có độ phản hồi cao, đồng thời đặt nền tảng cho các nghiên cứu tiếp theo về tối ưu hóa mạng lưới trạm sạc, phân tích hành vi người dùng và mở rộng kết nối phần cứng trong lĩnh vực giao thông điện. Các giải pháp kỹ thuật trong đề tài hoàn toàn có thể tham khảo và áp dụng cho các mô hình dịch vụ tương tự khác.


